% This section lists a number of 2-3 research questions that the thesis should address. 
% The research questions should be put into the context of the idea, and you should explain 
% what they mean and how you want to address each. We also suggest including an example scenario 
% of a successful outcome of the thesis, i.e., what would your tool allow us to do once it exists?

\chapter{Methodology}

\section{Proposed approach}

\subsubsection{Experimental platform}
As a subject of a research a simple robot will be taken, with two grippers and an elbow joint that can manipulate blocks (cuboctahedral voxels) one at a time.
An experimental platform will be realised as a simulation in Gazebo environment with control from ROS2 nodes. The  blocks and environment geometry robot joints and sensors and conections 
will be defined in Xacro/URDF files. 


\subsubsection{Generation of a model and specifications for formal-verification}
A custom program has to be written in Python, that parses the Xacro files and ROS2 Nodes controller to generate a NuSMV model:

  - a Boolean occupancy matrix representing the neighbourhood on a field relevant to the current replication.

  - replication program-counter plus local flags (gripper-occupied, block -held, etc.).

  The liveness and safety properties has to be defined independent from the simulation, in a form close to the following:
  \begin{itemize}
    
      \item Liveness. $G (start \rightarrow F (Robot_1==ready \land Robot_2==ready))$ — once replication starts, eventually both robots are able to function independently.

      \item Collision-free safety. $G \neg collision$ (or $AG \neg collision$ for CTL).

\end{itemize}


\section{Research tools}


Self-replication research in robotics involves complex interactions between 
assembly processes, modular design, environmental constraints, 
and control systems. Robotic simulation environments offer a valuable sandbox to 
explore these dynamics before physical implementation. 
The ideal simulation platform must support modularity, kinematics, 
sensing, and interaction with discrete components.

A set of tools for simulation and control of robotic systems including SRRS can be 
seen in the Table \ref{tab:tools}

\begin{table}[h!]
    \centering
    \begin{tabular}{|m{2cm}|m{1.5cm}|m{2cm}|m{2.5cm}|m{5cm}|}
        \hline
        \textbf{Simulator} & \textbf{Physics fidelity} & \textbf{Modularity support} &  \textbf{Custom control} & \textbf{Use case for SRRS} \\
        \hline
        Gazebo+ ROS &
        High &
        Strong via URDF/SDF &
        ROS integration, plugins, Python/C++ nodes &
        Physical replication with mobile robots and manipulators \\
        \hline
        CoppeliaSim (V-REP) &
        High &
        Native modular modeling &
        Lua, Python /remote APIs &
        Detailed control of part interaction and modular assembly \\
        \hline
        PyBullet &
        Medium to High &
        Good (via URDF /SDF) &
        Python API (fast iteration) &
        Cube/voxel-based systems and contact-rich interactions \\
        \hline
        Webots &
        Medium &
        Moderate &
        C/C++ /Python controllers &
        Lightweight simulation of small-scale replication setups \\
        \hline
        Unity+ \ \ ML-Agents &
        Medium &
        High (custom prefabs) &
        C\#,Python with \ \ \ \ \ ML-agents &
        Learning-based or perception-driven replication strategies \\
        \hline
        Isaac Sim (NVIDIA) &
        High &
        Strong (via Omniverse) &
        Python scripting, ROS2 &
        High-fidelity simulation of complex autonomous robots \\
        \hline
        Golly  \ \ \ NetLogo  CA tools &
        None &
        Tile-based, symbolic &
        Rule-based/ scripting &
        Theoretical studies (e.g. von Neumann or Langton automata) \\
        \hline
    \end{tabular}
    \caption{Comparison of simulation environments for self-replicating robotic systems}
    \label{tab:tools}
\end{table}

In order to simulate kinematic self-replicating systems and to be able to implement
control functions that are also applicable to a physical system, the Gazebo simulator 
with ROS infrastructure seems to be a plausible choice.

\begin{figure}[h]
    \centering
    \includegraphics[width=15cm]{Gazebo.png} 
    \caption{Gazebo Harmonic simulation interface}
    \label{fig:Gazebo}
\end{figure}

% \section{Analysis of results}


% \section{Evaluation}

